\chapter{项目背景与系统定位}

\section{赛题要求与业务背景}

ModelEngine 面向数据处理、知识生成和智能体应用提供开源工具链。DataMate 管理数据集、文件和算子任务，Nexent 管理智能体配置、工具绑定和对话执行。作品需要在两套平台之间建立可运行的医疗业务链路，完成数据处理、知识构建与问答、数据分析三个任务，并交付可核验的结果。

医疗资料同时包含自由文本、结构化表格、JSON 嵌套对象和 PDF 文档。实体、关系和来源信息分散在不同记录中，分析结果还需要以文字回答、统计图表和文件形式交付。

MediFlow 围绕三个赛题问题组织工程方案。数据处理阶段由智能体判断探查、处理和状态查询意图，服务根据数据状态选择格式处理链；知识构建阶段将记录转换为带来源的实体关系，经分级和校验后写入图谱；分析阶段把用户问题映射到图谱事实、分析表或只读查询，并生成可追踪的成果记录。

系统当前处理医学知识资料和聚合统计，定位为面向知识资料的工程系统。现有数据没有可用时间维度，分析结果按知识条目和关系事实解释。本文的系统对象包括三部分：MediFlow 提供 MCP 工具和医疗领域处理服务，Nexent 中的三个智能体负责任务级规划，DataMate 负责数据集和算子执行。


\section{MediFlow 系统定位}

MediFlow 是连接 Nexent 与 DataMate 的医疗数据处理与知识分析系统。系统把 DataMate 的数据集、算子和任务接口封装为 MCP 工具，再由领域服务补充格式路由、实体关系校验、图谱写入和只读查询规则。Nexent 中的三个智能体分别面向数据处理、知识图谱和数据分析；它们共享服务入口，并按职责限制可见工具和回复范围。

本文后续统一使用以下称谓：MediFlow 指本项目系统，MediFlow 服务指提供 MCP 接口的服务进程，智能体指 Nexent 中的任务配置，工具指智能体可调用的 MCP 接口，算子指 DataMate 执行的单步数据变换。系统输出中的“处理记录”保存文件状态、阶段统计和质量字段；“来源记录”保存数据集、文件和知识来源；“分析结果记录”保存问题、查询计划、SQL、数据行、图表配置和来源信息。

\section{数据驱动的业务链路}

作品沿“数据处理→知识构建→分析查询”的业务链路组织三个阶段。数据处理阶段接入原始资料，完成文本治理和结构保持后登记交付数据集；知识构建阶段读取交付数据集，抽取实体关系，经分级、校验和去重后写入知识图谱，并刷新分析库；分析查询阶段读取图谱和分析库，生成回答、图表和导出文件。
